\begin{frame}
	\frametitle{Détection automatique des évènements}
	\begin{block}{Idée}
		Dans cette slide, remettre en contexte la question de la détection des
		évènements. %
		Expliquer que cette première contribution propose une extension de l'analyse présentée dans la slide précédente.
	\end{block}
	\begin{itemize}
		\item Extension de l'analyse pour la détection d'évènements (création,
		      suppression, scission, \textit{etc.}),
		\item Enrichissement de la détection avec la distinction entre orbite de
		      suivi et d'origine. %
	\end{itemize}
\end{frame}

% \begin{frame}
% 	\frametitle{Création}
% 	\begin{block}{Idée}
% 		Règle de triangulation, présentation de la création de l'orbite sommet. %
% 		D'abord exprimer visuellement dans un objet, puis donner les conditions
% 		dans la règle. %
% 	\end{block}
% 	\begin{itemize}
% 		\item illustrer la création dans un objet,
% 		\item illustrer les conditions pour détecter la création d'une orbite
% 		      \begin{itemize}
% 			      \item uniquement des nœuds créés (à droite)
% 			      \item toutes les dimensions sont filtrées par l'ensemble des nœuds de
% 			            l'orbite
% 		      \end{itemize}
% 		\item introduire le terme d'orbite complète
% 	\end{itemize}
% \end{frame}

% \begin{frame}
% 	\frametitle{Suppression}
% 	\begin{block}{Idée}
% 		Expliquer que le problème de la suppression est symétrique à celui de la
% 		création. %
% 		C'est le cas dans les objets et dans les règles. %
% 		Ça devrait me permettre de ménager l'auditoire pour les évènements
% 		suivants (la scission et son symétrique~: la fusion). %
% 	\end{block}
% 	\begin{itemize}
% 		\item illustrer la suppression dans un objet,
% 		\item illustrer les conditions pour détecter la suppression d'une orbite,
% 		\item C'est le symétrique de la création.
% 	\end{itemize}
% \end{frame}

% \begin{frame}
% 	\frametitle{Scission}
% 	\begin{block}{Idée}
% 		Toucher un mot sur la notion de scission, expliquer que c'est la
% 		conséquence de toute division d'une orbite en plusieurs orbites de même
% 		type. %
% 		Ça comprend les subdivisions et les découpages d'orbites (par exemple
% 		l'insertion d'une arête dans une face ou l'intersection entre un volume et
% 		un plan). %
% 	\end{block}
% 	\begin{itemize}
% 		\item Illustrer la scission dans un objet,
% 		\item illustrer les conditions pour détecter la scission d'une orbite,
% 		\item Montrer qu'une scission est provoquée par la rupture de l'ensemble
% 		      des arcs implicites pour une dimension donnée. %
% 	\end{itemize}

% \end{frame}

% \begin{frame}
% 	\frametitle{Fusion}
% 	\begin{block}{Idée}
% 		Moment de pause pour l'auditoire, cet évènement est le symétrique de la scission. %
% 		Les conditions de détection sont aussi symétriques à ceux de la scission. %
% 	\end{block}
% 	\begin{itemize}
% 		\item Illustrer la fusion dans un objet,
% 		\item illustrer les conditions pour détecter la fusion de plusieurs
% 		      orbites,
% 		\item C'est le symétrique de la scission.
% 	\end{itemize}
% \end{frame}

% \begin{frame}
% 	\frametitle{Non-changement}
% 	\begin{block}{Idée}
% 		Présenter qu'une orbite filtrée qui n'est pas modifiée demeure inchangée
% 		dans la transformation. %
% 		Il faut garder cette slide simple et ne pas y passer trop de temps. %
% 	\end{block}
% 	\begin{itemize}
% 		\item Tout ce qui est filtré et inchangé dans une règle demeure inchangé dans l'objet,
% 		\item c'est aussi l'occasion de dire que ce qui n'est pas filtré ne peut
% 		      être transformé.
% 	\end{itemize}
% \end{frame}

% \begin{frame}
% 	\frametitle{Modification}
% 	\begin{block}{Idée}
% 		Une modification, au sens large, est tout ce qui n'est pas traité par les
% 		autres évènements. %
% 		Pour une face, c'est représenté par l'ajout ou la suppression d'une partie
% 		de ses sommets. %
% 	\end{block}
% 	\begin{itemize}
% 		\item Illustrer la modification dans un objet,
% 		\item ne pas s'étendre sur les conditions, c'est trop complexe,
% 		\item en pratique, on peut résumer ça a une augmentation ou réduction de
% 		      la topologie d'une cellule. %
% 	\end{itemize}

% \end{frame}

\begin{frame}
	\frametitle{Création}

	\begin{figure}
		\centering
		% \includegraphics<1>[height=.85\textheight]{evenement-creation-animation-grey}%
		% \includegraphics<2>[height=.85\textheight]{evenement-creation-animation-target}%
		% \includegraphics<3>[height=.85\textheight]{evenement-creation-animation-begin}%
		\includegraphics<1>[height=.85\textheight]{evenement-creation-animation-end}%
		\includegraphics<2>[height=.85\textheight]{evenement-creation-animation-end-2}%
	\end{figure}
\end{frame}

\begin{frame}{Scission d'une face $\orbit{0,1}$ par les liaisons implicites}
	\begin{figure}
		\centering
		% \includegraphics<1>[height=.85\textheight]{evenement-scission-impl-animation-grey}%
		% \includegraphics<2>[height=.85\textheight]{evenement-scission-impl-animation-target}%
		% \includegraphics<3>[height=.85\textheight]{evenement-scission-impl-animation-begin}%
		% \includegraphics<1>[height=.85\textheight]{evenement-scission-impl-animation-2-darts-impl-highlight}%
		\includegraphics<1>[height=.85\textheight]{evenement-scission-impl-animation-2-darts-impl-highlight-new}%
		\includegraphics<2>[height=.85\textheight]{evenement-scission-impl-animation-2-darts-impl-highlight-2}%
		\includegraphics<3>[height=.85\textheight]{evenement-scission-impl-animation-2-darts-impl-highlight-3}%
		\includegraphics<4>[height=.85\textheight]{evenement-scission-impl-animation-2-darts-impl-highlight-4}%
	\end{figure}
\end{frame}

\begin{frame}{Scission d'une arête $\orbit{0,2}$ par les liaisons explicites~: orbite complète}
		\centering
		% \includegraphics<1>[height=\textheight]{evenement-scission-expl-animation-grey.pdf}%
		% \includegraphics<2>[height=\textheight]{evenement-scission-expl-animation-target.pdf}%
		% \includegraphics<1>[height=\textheight]{evenement-scission-expl-animation-end.pdf}%
		\includegraphics<1>[height=.9\textheight]{evenement-scission-expl-animation-end-new.pdf}%
		\includegraphics<2>[height=.9\textheight]{evenement-scission-expl-animation-end-2.pdf}%
		\includegraphics<3>[height=.9\textheight]{evenement-scission-expl-animation-end-3.pdf}%
		\includegraphics<4>[height=.9\textheight]{evenement-scission-expl-animation-end-4.pdf}%
\end{frame}

\begin{frame}{Scission d'un sommet $\orbit{1,2}$ par les liaisons explicites~: orbite incomplète}
		\centering
		% \includegraphics<1>[height=\textheight]{false-positive-uncomplete-orbit-unsewed}%
		% \includegraphics<1>[height=\textheight]{false-positive-uncomplete-orbit-end}%
		\includegraphics<1>[height=.9\textheight]{false-positive-uncomplete-orbit-end-new}%
		\includegraphics<2>[height=.9\textheight]{false-positive-uncomplete-orbit-end-2}%
		\includegraphics<3>[height=.9\textheight]{false-positive-uncomplete-orbit-end-3}%
		\includegraphics<4>[height=.9\textheight]{false-positive-uncomplete-orbit-end-4}%
\end{frame}

\begin{frame}{Autres évènements}
	\begin{itemize}
		\item Suppression

		      \vspace{1em}

		\item Fusion

		      \vspace{1em}

		\item Non-changement

		      \vspace{1em}

		\item Autres modifications
	\end{itemize}
\end{frame}




%%% Local Variables:
%%% mode: LaTeX
%%% TeX-master: "../../main"
%%% End:
